\chapter{Το \textlatin{Dataset} και το Μοντέλο Δεδομένων}
\label{ch:dataset}

Το παρόν κεφάλαιο παρουσιάζει αναλυτικά το \textlatin{Plegma Dataset} που
χρησιμοποιήθηκε ως πηγή δεδομένων, καθώς
και τον τρόπο με τον οποίο τα δεδομένα αυτά αναπαρίστανται σε μορφή
\textlatin{NGSI-LD}. Ο στόχος είναι να γίνει κατανοητό, μέσα από πρακτικά
παραδείγματα, πώς το \textlatin{NGSI-LD} μοντελοποιεί πραγματικά δεδομένα
και πώς οι σχέσεις μεταξύ οντοτήτων επιτρέπουν πλούσια σημασιολογική
έκφραση.

% ------------------------------------------------------------------
\section{Το \textlatin{Plegma Dataset}}
\label{sec:plegma-dataset}
% ------------------------------------------------------------------

Το \textlatin{Plegma Dataset}~\cite{plegma2024} είναι μια ανοιχτή συλλογή
μετρήσεων κατανάλωσης ενέργειας και περιβαλλοντικών παραμέτρων από 13
νοικοκυριά στην Ελλάδα. Δημοσιεύτηκε από το \textlatin{University of
Strathclyde} και διατίθεται ελεύθερα για ερευνητικούς σκοπούς.

Το \textlatin{dataset} περιλαμβάνει δεδομένα που συλλέχθηκαν σε διάστημα
περίπου 15 μηνών (Ιούλιος 2022 με Σεπτέμβριος 2023) με χρονική ανάλυση ενός
λεπτού. Για κάθε νοικοκυριό διατίθενται:

\begin{itemize}
    \item Ηλεκτρολογικές μετρήσεις (\textlatin{Electric\_data}): Ολική
    ενεργός ισχύς (\textlatin{P\_agg} σε \textlatin{W}), τάση (\textlatin{V}),
    ρεύμα (\textlatin{A}) και ισχύς ανά συσκευή (κλιματιστικά, ψυγείο,
    πλυντήριο, θερμοσίφωνας).
    \item Περιβαλλοντικές μετρήσεις (\textlatin{Environmental\_data}):
    Εσωτερική και εξωτερική θερμοκρασία (°\textlatin{C}) και υγρασία (\%).
    \item Μεταδεδομένα κτιρίου (\textlatin{appliances\_metadata.csv},
    \textlatin{houses\_metadata.xlsx}): Χαρακτηριστικά κτιρίου, σύστημα θέρμανσης,
    έτος κατασκευής, αριθμός δωματίων, κατάσταση ένοικου.
    \item Δημογραφικά στοιχεία: Ανωνυμοποιημένα σωκοδημογραφικά
    δεδομένα του νοικοκύρη (ηλικιακή ομάδα, εκπαιδευτικό επίπεδο, εισόδημα).
\end{itemize}

Για 13 νοικοκυριά, μετρήσεις ανά λεπτό, σε διάστημα 15 μηνών, ο συνολικός
όγκος ανέρχεται σε περίπου 8,1 εκατομμύρια χρονοσφραγίδες ανά τύπο μέτρησης.
Εάν προσπαθούσαμε να φορτώσουμε όλα αυτά τα δεδομένα ως πάροχος, αποδεικνύοντας κάθε φορά την
ταυτότητά μας θα θέλαμε περίπου 1 ημέρα για ολοκληρωθεί η διαδικασία. Έτσι τα δεδομένα 
φορτώθηκα στο σύστημα παρακάμπτοντας την διαδικασία πιστοποίησης.

% ------------------------------------------------------------------
\section{Οντότητα \textlatin{Building}}
\label{sec:entity-building}
% ------------------------------------------------------------------

Η οντότητα \textlatin{Building} αναπαριστά ένα νοικοκυριό και αποθηκεύει
τα στατικά μεταδεδομένα του. Εμπνέεται από το επίσημο μοντέλο
\textlatin{smart-data-models/dataModel.Building}. Έτσι έχει την μορφή:

\vspace{0.5cm}
\selectlanguage{english}
\begin{lstlisting}[language={}, caption={Building entity for House01 in NGSI-LD.}]
{
  "id": "urn:ngsi-ld:Building:House01",
  "type": "https://plegma.example.org/vocab#Building",

  "https://smartdatamodels.org/dataModel.Weather/category": {
    "type": "Property",
    "value": ["residential", "apartment"]
  },
  "dateOfConstruction": {
    "type": "Property",
    "value": "1950/1970"
  },
  "numberOfRooms": {
    "type": "Property",
    "value": 3
  },
  "https://plegma.example.org/vocab#occupancyStatus": {
    "type": "Property",
    "value": "renter"
  },
  "https://plegma.example.org/vocab#heatingSystem": {
    "type": "Property",
    "value": ["radiatorOil", "airConditioner"]
  },
  "https://plegma.example.org/vocab#waterHeaterType": {
    "type": "Property",
    "value": "electricBoiler"
  },
  "https://plegma.example.org/vocab#hasSolarPanels": {
    "type": "Property",
    "value": false
  },
  "https://plegma.example.org/vocab#hasOccupant": {
    "type": "Relationship",
    "object": "urn:ngsi-ld:Person:House01:Occupant"
  },
  "https://plegma.example.org/vocab#hasAppliance": {
    "type": "Relationship",
    "object": [
      "urn:ngsi-ld:Device:House01:ac_1",
      "urn:ngsi-ld:Device:House01:ac_2",
      "urn:ngsi-ld:Device:House01:boiler",
      "urn:ngsi-ld:Device:House01:fridge",
      "urn:ngsi-ld:Device:House01:washing_machine"
    ]
  },
  "@context": "https://example.org/plegma-context.jsonld"
}
\end{lstlisting}
\selectlanguage{greek}

Αξίζει να παρατηρηθούν τα εξής:
\begin{itemize}
    \item Η ιδιότητα \textlatin{hasAppliance} είναι Σχέση
    (\textlatin{Relationship}) και δείχνει σε ένα σύνολο από
    \textlatin{Device URI}. Αυτό επιτρέπει ερωτήματα όπως «ποιες συσκευές
    ανήκουν σε αυτό το κτίριο» χωρίς να χρειαστεί \textlatin{JOIN} ή
    εξωτερική αναζήτηση.
    \item Οι ιδιότητες που ανήκουν στο \textlatin{Plegma} προσαρτώνται με
    το \textlatin{namespace} \textlatin{https://plegma.example.org/vocab\#},
    ενώ οι κοινές ιδιότητες χρησιμοποιούν τον
    κοινό \textlatin{namespace} \textlatin{smart-data-models}.
    \item Το αναγνωριστικό (\textlatin{id}) ακολουθεί το πρότυπο
    \textlatin{urn:ngsi-ld:\{EntityType\}:\{HouseXX\}}.
\end{itemize}

% ------------------------------------------------------------------
\section{Οντότητα \textlatin{Device}}
\label{sec:entity-device}
% ------------------------------------------------------------------

Η οντότητα \textlatin{Device} αναπαριστά μια μετρούμενη οικιακή συσκευή.
Δημιουργούνται πέντε \textlatin{Device} οντότητες ανά νοικοκυριό, μία για κάθε
συσκευή που καταγράφεται στο \textlatin{dataset} (δύο κλιματιστικά,
θερμοσίφωνας, ψυγείο, πλυντήριο).

\vspace{0.5cm}
\selectlanguage{english}
\begin{lstlisting}[language={}, caption={Device entity for air conditioner ac\_1.}]
{
  "id": "urn:ngsi-ld:Device:House01:ac_1",
  "type": "https://plegma.example.org/vocab#Device",

  "category": {
    "type": "Property",
    "value": ["hvac"]
  },
  "controlledProperty": {
    "type": "Property",
    "value": ["power"]
  },
  "https://plegma.example.org/vocab#applianceType": {
    "type": "Property",
    "value": "airConditioner"
  },
  "https://plegma.example.org/vocab#ratedWattage": {
    "type": "Property",
    "value": 1800,
    "unitCode": "WTT"
  },
  "https://plegma.example.org/vocab#detectionThreshold": {
    "type": "Property",
    "value": 50,
    "unitCode": "WTT"
  },
  "https://plegma.example.org/vocab#installedIn": {
    "type": "Relationship",
    "object": "urn:ngsi-ld:Building:House01"
  },
  "@context": "https://example.org/plegma-context.jsonld"
}
\end{lstlisting}
\selectlanguage{greek}

% ------------------------------------------------------------------
\section{Οντότητα \textlatin{HouseholdElectricMeasurement}}
\label{sec:entity-electric}
% ------------------------------------------------------------------

Η οντότητα \textlatin{HouseholdElectricMeasurement} αναπαριστά ένα στιγμιότυπο
ηλεκτρολογικών μετρήσεων από τον έξυπνο μετρητή ενός νοικοκυριού. Είναι η
πιο πυκνή οντότητα του \textlatin{dataset}.

\vspace{0.5cm}
\selectlanguage{english}
\begin{lstlisting}[language={}, caption={HouseholdElectricMeasurement entity for House01.}]
{
  "id": "urn:ngsi-ld:HouseholdElectricMeasurement:House01:2022-07-15T00:00:00Z",
  "type": "https://plegma.example.org/vocab#HouseholdElectricMeasurement",

  "https://smartdatamodels.org/dateObserved": {
    "type": "Property",
    "value": { "@type": "DateTime", "@value": "2022-07-15T00:00:00Z" }
  },
  "https://plegma.example.org/vocab#activePower": {
    "type": "Property",
    "value": 121.542,
    "unitCode": "WTT",
    "observedAt": "2022-07-15T00:00:00.000Z"
  },
  "https://plegma.example.org/vocab#phaseVoltage": {
    "type": "Property",
    "value": 236.83,
    "unitCode": "VLT",
    "observedAt": "2022-07-15T00:00:00.000Z"
  },
  "https://plegma.example.org/vocab#current": {
    "type": "Property",
    "value": 0.624,
    "unitCode": "AMP",
    "observedAt": "2022-07-15T00:00:00.000Z"
  },
  "https://plegma.example.org/vocab#appliancePower": {
    "type": "Property",
    "value": {
      "ac_1": 0,
      "ac_2": 2.033,
      "boiler": 0,
      "fridge": 45.2,
      "washing_machine": 0
    },
    "unitCode": "WTT",
    "observedAt": "2022-07-15T00:00:00.000Z"
  },
  "https://plegma.example.org/vocab#dataQualityIssue": {
    "type": "Property",
    "value": false
  },
  "https://plegma.example.org/vocab#refBuilding": {
    "type": "Relationship",
    "object": "urn:ngsi-ld:Building:House01"
  },
  "@context": "https://example.org/plegma-context.jsonld"
}
\end{lstlisting}
\selectlanguage{greek}

Η ισχύς ανά συσκευή αποθηκεύεται ως ένα αντικείμενο
\textlatin{appliancePower} με τιμή λεξικού (π.χ. \textlatin{\{"ac\_1": 0, "fridge":
45.2, ...\}}), αντί να δημιουργείται ξεχωριστή οντότητα για κάθε συσκευή ανά
χρονοσφραγίδα. Η επιλογή αυτή βασίστηκε στο ότι η δημιουργία ξεχωριστών οντοτήτων ανά συσκευή
ανά λεπτό θα αύξανε τον αριθμό οντοτήτων κατά 5 φορές ($\approx$3,4
εκατομμύρια αντί για $\approx$560 χιλιάδες ανά μήνα για 13 νοικοκυριά),
χωρίς ουσιαστικό πλεονέκτημα για τα ερωτήματα του παρόντος
\textlatin{Data Space}. Επίση, το κλειδί του λεξικού (π.χ.
\textlatin{ac\_1}) αντιστοιχεί άμεσα στο επίθημα του \textlatin{URI} της
αντίστοιχης \textlatin{Device} οντότητας
(\textlatin{urn:ngsi-ld:Device:House01:ac\_1}), διατηρώντας τη σημασιολογική
σύνδεση χωρίς επιπλέον οντότητες.


% ------------------------------------------------------------------
\section{Οντότητα \textlatin{EnvironmentalMeasurement}}
\label{sec:entity-environmental}
% ------------------------------------------------------------------

Η οντότητα \textlatin{EnvironmentalMeasurement} αναπαριστά ένα στιγμιότυπο
περιβαλλοντικών μετρήσεων. Εμπνέεται από το μοντέλο
\textlatin{smart-data-models/dataModel.Weather/WeatherObserved} αλλά διακρίνει
μεταξύ εσωτερικών και εξωτερικών μετρήσεων.

\vspace{0.5cm}
\selectlanguage{english}
\begin{lstlisting}[language={}, caption={EnvironmentalMeasurement entity for House01.}]
{
  "id": "urn:ngsi-ld:EnvironmentalMeasurement:House01:2022-07-15T00:00:00Z",
  "type": "https://plegma.example.org/vocab#EnvironmentalMeasurement",

  "https://smartdatamodels.org/dateObserved": {
    "type": "Property",
    "value": { "@type": "DateTime", "@value": "2022-07-15T00:00:00Z" }
  },
  "https://plegma.example.org/vocab#indoorTemperature": {
    "type": "Property",
    "value": 27.4,
    "unitCode": "CEL",
    "observedAt": "2022-07-15T00:00:00.000Z"
  },
  "https://plegma.example.org/vocab#indoorHumidity": {
    "type": "Property",
    "value": 32,
    "unitCode": "P1",
    "observedAt": "2022-07-15T00:00:00.000Z"
  },
  "https://plegma.example.org/vocab#outdoorTemperature": {
    "type": "Property",
    "value": 34.1,
    "unitCode": "CEL",
    "observedAt": "2022-07-15T00:00:00.000Z"
  },
  "https://plegma.example.org/vocab#outdoorHumidity": {
    "type": "Property",
    "value": 40,
    "unitCode": "P1",
    "observedAt": "2022-07-15T00:00:00.000Z"
  },
  "https://plegma.example.org/vocab#refBuilding": {
    "type": "Relationship",
    "object": "urn:ngsi-ld:Building:House01"
  },
  "@context": "https://example.org/plegma-context.jsonld"
}
\end{lstlisting}
\selectlanguage{greek}

% ------------------------------------------------------------------
\section{Το \textlatin{JSON-LD Context} του \textlatin{Plegma Dataset}}
\label{sec:plegma-context}
% ------------------------------------------------------------------

Ένα κρίσιμο στοιχείο του \textlatin{NGSI-LD} είναι το \textlatin{context}
(\textlatin{@context}), το οποίο αντιστοιχίζει τα σύντομα ονόματα των
ιδιοτήτων σε παγκόσμια \textlatin{URI}. Αυτό επιτρέπει στους αναγνώστες
να κατανοούν τη σημασία κάθε ιδιότητας ανεξάρτητα από τον οργανισμό που τη
δημοσίευσε.

Το αρχείο \textlatin{plegma-context.jsonld} ορίζει τρία επίπεδα:

\begin{enumerate}
    \item Εισαγωγή επίσημων \textlatin{context}: Ο \textlatin{Plegma
    context} εισάγει πρώτα τα επίσημα \textlatin{context} της
    \textlatin{ETSI} (\textlatin{NGSI-LD core}) και των
    \textlatin{smart-data-models}, κληρονομώντας τον ορισμό κοινών
    ιδιοτήτων όπως \textlatin{dateObserved}, \textlatin{location} κ.λπ.

    \item \textlatin{Namespace} \textlatin{plegma:}: Ορίζεται μια
    σύντομη αναφορά στο \textlatin{namespace} του \textlatin{Plegma}:

\selectlanguage{english}
\begin{lstlisting}[language={}]
"plegma": "https://plegma.example.org/vocab#"
\end{lstlisting}
\selectlanguage{greek}

    \item Αντιστοίχιση Ιδιοτήτων: Κάθε ιδιότητα που ορίζεται
    στα \textlatin{schemas} αντιστοιχίζεται σε ένα \textlatin{URI}:
\selectlanguage{english}
\begin{lstlisting}[language={}, caption={Context mapping excerpt.}]
{
  "@context": {
    "plegma": "https://plegma.example.org/vocab#",
    "Building":  "plegma:Building",
    "Device":    "plegma:Device",
    "HouseholdElectricMeasurement": "plegma:HouseholdElectricMeasurement",
    "EnvironmentalMeasurement": "plegma:EnvironmentalMeasurement",
    "activePower":    { "@id": "plegma:activePower",    "@type": "@id" },
    "phaseVoltage":   { "@id": "plegma:phaseVoltage",   "@type": "@id" },
    "appliancePower": { "@id": "plegma:appliancePower"  },
    "indoorTemperature":  { "@id": "plegma:indoorTemperature"  },
    "indoorHumidity":     { "@id": "plegma:indoorHumidity"     },
    "outdoorTemperature": { "@id": "plegma:outdoorTemperature" },
    "outdoorHumidity":    { "@id": "plegma:outdoorHumidity"    },
    "refBuilding":    { "@id": "plegma:refBuilding", "@type": "@id" },
    "hasOccupant":    { "@id": "plegma:hasOccupant", "@type": "@id" },
    "hasAppliance":   { "@id": "plegma:hasAppliance", "@type": "@id" }
  }
}
\end{lstlisting}
\selectlanguage{greek}
\end{enumerate}

Ο \textlatin{Orion-LD} απαιτεί το \textlatin{context} να είναι προσβάσιμο
μέσω \textlatin{HTTP} κατά τη δημιουργία οντοτήτων. Αυτό επιτρέπει στον
\textlatin{broker} να επιλύσει τα ονόματα ιδιοτήτων στα πλήρη
\textlatin{URI} τους εσωτερικά. Στη διάρκεια της επίδειξης, χρησιμοποιήθηκε
το επίσημο \textlatin{context} των \textlatin{smart-data-models}:

\selectlanguage{english}
\begin{lstlisting}[language={}]
CONTEXT_URL = "https://raw.githubusercontent.com/smart-data-models/
               dataModel.Energy/master/context.jsonld"
\end{lstlisting}
\selectlanguage{greek}

% ------------------------------------------------------------------
\section{Αντιστοίχιση \textlatin{CSV} σε \textlatin{NGSI-LD}}
\label{sec:csv-mapping}
% ------------------------------------------------------------------

Οι παρακάτω πίνακες  παρουσιάζουν την αντιστοίχιση μεταξύ
στηλών \textlatin{CSV} και ιδιοτήτων \textlatin{NGSI-LD} για τα ηλεκτρολογικά
δεδομένα.

\begin{table}[H]
\centering
\caption{Αντιστοίχιση στηλών \textlatin{CSV} σε ιδιότητες \textlatin{NGSI-LD} (ηλεκτρολογικά δεδομένα).}
\label{tab:electric-mapping}
\begin{tabular}{|>{\centering\arraybackslash}m{3.5cm}|>{\centering\arraybackslash}m{5cm}|>{\centering\arraybackslash}m{2cm}|>{\centering\arraybackslash}m{2cm}|}
\hline
\textbf{Στήλη \textlatin{CSV}} & \textbf{Ιδιότητα \textlatin{NGSI-LD}} & \textbf{Μονάδα} & \textbf{Τύπος} \\
\hline
\textlatin{timestamp} & \textlatin{dateObserved} & - & \textlatin{Property} \\
\hline
\textlatin{P\_agg} & \textlatin{plegma:activePower} & \textlatin{W} & \textlatin{Property} \\
\hline
\textlatin{V} & \textlatin{plegma:phaseVoltage} & \textlatin{V} & \textlatin{Property} \\
\hline
\textlatin{A} & \textlatin{plegma:current} & \textlatin{A} & \textlatin{Property} \\
\hline
\textlatin{ac\_1} & \textlatin{plegma:appliancePower.ac\_1} & \textlatin{W} & \textlatin{Property} \\
\hline
\textlatin{ac\_2} & \textlatin{plegma:appliancePower.ac\_2} & \textlatin{W} & \textlatin{Property} \\
\hline
\textlatin{boiler} & \textlatin{plegma:appliancePower.boiler} & \textlatin{W} & \textlatin{Property} \\
\hline
\textlatin{fridge} & \textlatin{plegma:appliancePower.fridge} & \textlatin{W} & \textlatin{Property} \\
\hline
\textlatin{washing\_machine} & \textlatin{washing\_machine} & \textlatin{W} & \textlatin{Property} \\
\hline
\textlatin{issues} & \textlatin{plegma:dataQualityIssue} & - & \textlatin{Property (bool)} \\
\hline
- & \textlatin{plegma:refBuilding} & - & \textlatin{Relationship} \\
\hline
\end{tabular}
\end{table}

\begin{table}[H]
\centering
\caption{Αντιστοίχιση στηλών \textlatin{CSV} σε ιδιότητες \textlatin{NGSI-LD} (περιβαλλοντικά δεδομένα).}
\label{tab:env-mapping}
\begin{tabular}{|>{\centering\arraybackslash}m{4cm}|>{\centering\arraybackslash}m{5cm}|>{\centering\arraybackslash}m{2cm}|>{\centering\arraybackslash}m{2cm}|}
\hline
\textbf{Στήλη \textlatin{CSV}} & \textbf{Ιδιότητα \textlatin{NGSI-LD}} & \textbf{Μονάδα} & \textbf{Τύπος} \\
\hline
\textlatin{timestamp} & \textlatin{dateObserved} & - & \textlatin{Property} \\
\hline
\textlatin{internal\_temperature} & \textlatin{plegma:indoorTemperature} & °\textlatin{C} & \textlatin{Property} \\
\hline
\textlatin{internal\_humidity} & \textlatin{plegma:indoorHumidity} & \% & \textlatin{Property} \\
\hline
\textlatin{external\_temperature} & \textlatin{plegma:outdoorTemperature} & °\textlatin{C} & \textlatin{Property} \\
\hline
\textlatin{external\_humidity} & \textlatin{plegma:outdoorHumidity} & \% & \textlatin{Property} \\
\hline
\end{tabular}
\end{table}